% Bagian: second-order-logic
% Bab: metatheory
% Subbagian: introduction

\documentclass[../../../include/open-logic-section]{subfiles}

\begin{document}

\olfileid[id]{sol}{met}{int}

\olsection{Pendahuluan}

Logika orde pertama mempunyai sejumlah sifat yang baik. Kita tahu bahwa
logika ini tidak dapat diputuskan, tetapi setidaknya dapat diaksiomatisasi.
Artinya, terdapat sistem pembuktian bagi logika orde pertama yang sahih dan
lengkap, yakni sistem tersebut menghasilkan relasi !!a{derivability}~$\Proves$
dengan sifat bahwa untuk setiap himpunan !!{sentence}~$\Gamma$ dan setiap
!!{sentence}~$!A$, $\Gamma \Entails !A$ jika dan hanya jika $\Gamma \Proves
!A$. Secara khusus, ini berarti bahwa himpunan validitas logika orde pertama
!!{computably enumerable}. Terdapat suatu fungsi yang dapat dihitung~$f\colon
\Nat \to \Sent[L]$ sedemikian sehingga nilai-nilai~$f$ tepat merupakan semua
!!{sentence} valid dari~$\Lang{L}$. Hal ini berlaku karena !!{derivation}s dapat
dienumerasi, lalu derivasi-derivasi yang !!{derive} satu~!!{sentence} dipetakan
ke !!{sentence} tersebut. Logika orde kedua lebih ekspresif daripada logika
orde pertama, sehingga secara umum lebih rumit untuk mencakup semua
validitasnya. Bahkan, kita akan menunjukkan bahwa logika orde kedua bukan
hanya tidak dapat diputuskan, melainkan validitas-validitasnya bahkan tidak
!!{computably enumerable}. Ini berarti bahwa tidak mungkin ada sistem
pembuktian yang sahih dan lengkap bagi logika orde kedua (meskipun tersedia
sistem pembuktian yang sahih tetapi tidak lengkap, dan sistem-sistem tersebut
bahkan merupakan objek penelitian yang penting).

Logika orde pertama juga mempunyai dua sifat lain: logika ini kompak (jika
setiap himpunan bagian berhingga dari suatu himpunan !!{sentence}~$\Gamma$
dapat dipenuhi, maka $\Gamma$ sendiri dapat dipenuhi) dan Teorema
L\"owenheim--Skolem berlaku baginya (jika $\Gamma$ mempunyai model tak hingga,
maka himpunan itu mempunyai model !!{denumerable}). Kedua hasil ini tidak berlaku
bagi logika orde kedua. Sekali lagi, alasannya ialah bahwa logika orde kedua
dapat mengungkapkan fakta-fakta tentang ukuran !!{domain} yang tidak dapat
diungkapkan oleh logika orde pertama.

\end{document}
